\documentclass[a4paper,12pt]{article}

\usepackage[a4paper,margin=16mm]{geometry}
\usepackage[german,main=french]{babel}
\usepackage{fontspec}
\usepackage{microtype}

\setmainfont{EBGaramond}
[
  Extension      = .otf,
  UprightFont    = *-Regular,
  ItalicFont     = *-Italic,
  BoldFont       = *-Bold,
  BoldItalicFont = *-BoldItalic,
  Numbers        = Lowercase
]

\usepackage{paracol}
\newcommand\chunks[2]{%
\begin{leftcolumn*}\begin{otherlanguage}{german}%
{#1}%
\end{otherlanguage}\end{leftcolumn*}%
\begin{rightcolumn}\begin{otherlanguage}{french}%
{#2}%
\end{otherlanguage}\end{rightcolumn}%
}

\newcommand\strophe[1]{\subsection*{\normalsize\normalfont\itshape #1}}

\begin{document}\thispagestyle{empty}
\setlength{\columnsep}{8mm}


\begin{paracol}{2}

  \chunks{\section*{Ach, ich fühl’s, es ist verschwunden}}
         {\section*{Ah je le sens, il a disparu}}

  \chunks{\noindent\raggedright\small Wolfgang Amadeus Mozart, \textit{Die Zauberflöte}}
         {\noindent\raggedright\small Wolfgang Amadeus Mozart, \textit{La Flûte enchantée}}

  \chunks{\noindent\raggedright\small Libretto von Emanuel Schikaneder (1791)}
         {\noindent\raggedright\small Livret d’Emanuel Schikaneder (1791)}

  \chunks{}
         {\noindent\raggedright\small Texte français : Jean Spenlehauer}

  \chunks{\strophe{Zweiter Aufzug, achtzehnter Auftritt, Nr.~17, Arie}}
         {\strophe{Deuxième acte, dix-huitième scène, n\textsuperscript{o}~17, air}}

  \chunks{\begin{verse}
          \textsc{Pamina}\\[\medskipamount]
          Ach ich fühl’s, es ist verschwunden!\\
          Ewig hin der Liebe Glück!\\
          Nimmer kommt ihr Wonnestunden,\\
          Meinem Herzen mehr zurück.
          \end{verse}}
         {\begin{verse}
          \textsc{Pamina}\\[\medskipamount]
          Ah je le sens, il a disparu\\
          Pour toujours le bonheur de l’amour !\\
          Plus jamais, heures délicieuses\\
          Vous ne reviendrez à mon cœur.
          \end{verse}}

  \chunks{\begin{verse}
          Sieh Tamino, diese Tränen\\
          Fliessen Trauter, dir allein.\\
          Fühlst du nicht der Liebe Sehnen,\\
          So wird Ruh, so wird Ruh im Tode sein.
          \end{verse}}
         {\begin{verse}
          Vois, Tamino, ces larmes\\
          Coulent pour toi seul, bien-aimé.\\
          Si tu ne ressens pas le désir de l’amour,\\
          Alors ce sera la paix, ce sera la paix dans la mort.
          \end{verse}}
\end{paracol}
\end{document}
